% 本文件由 LaTeX 排版 Agent 根据有效 translation.json 自动生成。
% author=John Chrysostom; work_id=1907; title=论谦卑之心

\chapter{论谦卑之心}

% structure: p0001 source=h1 target=omitted reason=page-title-already-rendered-by-parent

本篇讲道开始时，暗指不久前关于法利塞人与税吏比喻的评论。这些话出现在\Person{金口圣若望}的第五篇\WorkTitle{反阿诺米派讲道}中，这是一组讲道中的一篇，根据内部证据，可推定其写作年份为386年底或387年初。因此，以下这篇讲道是在\Place{安提约基雅}宣讲的，很可能是在386年圣诞节前夕。当时有些人解释标题所引用的\Person{圣保禄}的话，认为只要宣讲了基督，所教导的实际教义是正确还是异端都无关紧要。本篇讲道的主要目的，就是要为宗徒的话辩护，驳斥这种错误而有害的解释。\par

1. 当我们近日提及法利塞人与税吏，并比喻性地用美德与恶习驾起两辆战车时，我们指出了两方面的真理：谦卑之心的获益何其大，而骄傲的损害又何其大。因为骄傲即使与正义、禁食和什一奉献相结合，仍落在后面；而谦卑之心即使与罪恶相联，也超越了法利塞人的那对马车，尽管它的御者并不优秀。有什么比税吏更糟糕的呢？但同样，由于他使自己的灵魂痛悔，并称自己为罪人——而这正是他的真实状况——他胜过了那位法利塞人，后者既有禁食可夸，又有什一奉献可提，并且远离任何恶行。这是为何？通过什么？因为即使法利塞人远离了贪婪和抢劫，他却将万恶之母——虚荣与骄傲——植根于自己的灵魂。因此\Person{保禄}也劝勉说：\Quote{各人应当考验自己的行为，这样，他夸耀只在自己身上，而不在别人身上。}而那法利塞人却公开站出来控告全世界，声称自己比所有活着的人更好。即便他仅仅把自己放在十个人、五个人、两个人或一个人之前，这也是不可容忍的；但他不仅把自己置于全世界之前，还控告所有人。因此他在赛跑中落后了。正如一艘船经历了无数波涛，逃过了许多风暴，却在进港处撞上礁石，失去了舱内所有的宝藏——同样，这位法利塞人经历了禁食及其他一切美德的辛劳之后，因未能管住自己的舌头，在港口本身遭受了货物沉没的厄运。因为从祈祷中回家，他本应获得收益，却遭受了如此巨大的损失，这无异于在港口遇难。\par

2. 因此，亲爱的各位，即使我们已经攀登到美德的最顶峰，也要将自己视为最末的；因为我们知道，骄傲能将不加提防的人从天上拉下来，而谦卑之心能将知道保持清醒的人从罪恶的深渊高举到天上。正是这一点使税吏位于法利塞人之前；而骄傲，我指的是傲慢自大的精神，甚至胜过了一种无形的势力，即\Person{魔鬼}的势力；而谦卑之心及承认自己所犯的罪恶，却使强盗在宗徒之前进入了乐园。如果那些承认自己罪恶的人为自己带来的信心如此巨大，那么那些自觉有许多善行却谦卑自己灵魂的人，将赢得何等大的冠冕呢？因为当罪恶与谦卑之心结合时，它以如此轻盈的速度奔跑，以至于超越并胜过了与骄傲结合的正义。因此，如果你将谦卑之心与正义结合，它将到达何处？它将穿过多少层天？它必定会在天主的宝座前停下步伐，在天使中间充满信心。另一方面，如果骄傲与正义结合，以其罪恶的过重和分量足以拖垮\Emph{它的}信心；如果它与罪恶结合，它将能把具有它的人推入何等深的地狱呢？我说这些话，不是要我们忽视正义，而是要我们避免骄傲；不是要我们犯罪，而是要我们保持清醒。因为谦卑之心是我们所拥有的爱智之德的基础。即使你已经建造了无数的事物：即使有施舍、祈祷、禁食和一切美德；除非先奠定这个基础，否则一切建造都将徒劳无益，并且会像建在沙土上的房屋一样轻易倒塌。因为没有任何一样善行不需要它；没有它，任何善行都无法站立。即使你提到节制、童贞、轻视钱财或任何其他事物，如果缺乏谦卑之心，一切都是不洁、可憎和令人厌恶的。因此，让我们在一切事上带着她——在言语、行为、思想上，并以此为基础建立这些恩宠。\par

3. 关于谦卑心意的事，已经说得足够多了；不是因为这德行的价值（因为没有人能按它的价值来颂扬它），而是为了你们爱德的明悟。因为我深知，即使从所说的少许话中，你们也会以极大的热忱去拥抱它。但由于也有必要澄清并显明今天所诵读的宗徒言论，这言论在许多人看来似乎为懈怠提供了借口，免得有些人从中为自己预备了一种冰冷的辩护，而疏忽了自己的救恩——让我们将讲论转向这一点。那么，这言论是什么呢？经上说：\Quote{或出于假意，或出于诚心，基督总被宣讲。}\BibleRef{斐 1:18}许多人完全随意地曲解这话，不阅读前后的经文；而是将其从其余部分的连贯中割裂出来，为了自己灵魂的丧亡，将它提供给较为懈怠的人。因为当他们试图引诱人离开纯正的信仰时，看见他们害怕战栗，并以这样做并非没有危险为由，想要缓解他们的恐惧，于是提出这宗徒的宣告，说：\Quote{或出于假意或出于诚心，让基督被传扬吧。}但这些事并非如此，绝非如此。首先，他没有说\Quote{让他被传扬}，而是\Quote{他被传扬}，这两者之间的区别是很大的。因为\Quote{让他被传扬}这句话属于立法者；而\Quote{他被传扬}则属于宣告事件的人。因为保禄并非制定一条准许异端存在的法律，而是劝诫所有听他话的人，请听他所说的：\Quote{若有人给你们宣讲的福音，与你们所接受的不同，就让他当受咒诅，即便是我，或是来自天上的天使。}\BibleRef{迦 1:8}如果他知道这事没有危险，他绝不会咒诅自己和天使。再者——\Quote{我以天主的妒爱妒爱你们，}他说，\Quote{因为我曾把你们许配给一个丈夫，一个贞洁的童女；我害怕，恐怕正如蛇用狡猾欺骗了厄娃，你们的心意也会从向基督所有的纯一受到败坏。}请看，他既设定了纯一，又不容任何例外。因为若有例外，就没有危险；若没有危险，保禄就不会害怕；基督也不会命令将莠子烧掉，如果听从这个或那个或另一个，或一律听从，本是无关紧要的事。\par

4. 那么，究竟是什么意思？我愿意从稍早之处向你们叙述全部历史；因为必须知道\Person{保禄}在写信时身处何种境况。那么，他当时身处何种境况？在监狱、锁链和难以忍受的危险中。这一切从何处显明？从书信本身。因为在此以前他说：\BibleQuote{弟兄们，我愿意你们知道，我的处境反而更促进了福音的进展，以致我的锁链在基督内，在整个宫廷和其余的人中，都成了明显的；并且许多弟兄因我的锁链而信赖主，更加大胆地毫无畏惧地宣讲圣言。}\BibleRef{斐 1:12} 那时\Person{尼禄}将他投入监狱。就如某个强盗进入房屋，当众人沉睡时，偷窃一切；若他看见有人点了灯，便将灯熄灭，又将持灯者杀死，以便他能安全地偷盗和掠夺他人的财物；同样，凯撒\Person{尼禄}那时也如同强盗和窃贼，当众人都陷入深沉无知觉的睡眠时，掠夺所有人的财物，闯入卧室，颠覆房屋，显示各种邪恶；当他看见\Person{保禄}在全世点了一盏灯（即他教导的圣言），并指责他的邪恶时，便竭力要扑灭所宣讲的，并将教师除掉，以便他能有权势地随意行事；于是将那位圣人捆绑，投入监狱。正是在那时，真福\Person{保禄}写了这些事。谁不会震惊？谁不会赞叹？或者，谁能足够地惊叹和钦慕那高贵而达于天堂的灵魂？他身在罗马被捆绑囚禁，在如此遥远的距离，写信给斐理伯人？因为你们知道马其顿和罗马之间的距离有多远。但是，路途的遥远、所需的时间、事务的繁忙、接踵而来的危险和患难，没有什么能驱散他对门徒们的爱和记忆；他心中保留着他们所有人；他的双手被锁链捆绑得如此紧，但他的灵魂却因对门徒的渴望而更加紧密地联结和钉牢。他自己在书信的序言中也宣告说：\BibleQuote{因为我在心内怀有你们，无论在我的锁链中，或在福音的辩护和证实中。}\BibleRef{斐 1:7} 正如一位君王，清晨登上宝座，坐在王宫中，立即从四面八方收到无数的信件；同样，他也如同坐在王宫中，坐在地牢里，收到和发出的信件远为更多；各民族从各处将他智慧中关于自身所发生的一切都报告给他；他处理的事务比在位的君主更多，因为托付给他的管辖区更大。因为事实上，天主不仅把罗马人居住的地区交在他手里，也把一切外邦人，无论陆地海洋。为表明这一点，他对罗马人说：\BibleQuote{弟兄们，我不愿意你们不知道，我多次决意到你们那里去，但一直受到阻碍，直到现在；为要在你们中间也得些果实，如同在其他外邦人中一样。无论对希腊人、对化外人、对智者、对愚者，我都欠债。}\BibleRef{罗 1:13} 因此，他每天焦虑，一时为格林多人，一时为马其顿人；斐理伯人如何，卡帕多细亚人如何，迦拉达人如何，雅典人如何，居住在本都的人如何，所有人整体如何。但尽管如此，他把整个世界都交在手中，他不仅不断关怀整个民族，也关怀每一个人；他时而为\Person{敖尼息摩}寄出信函，时而为格林多中犯奸淫的人。因为他从不这样看——那是一个犯了罪需要辩护的个体；而是一个人；一个人，天主最宝贵的生灵；且为了他，圣父连独生子都没有顾惜。\par

5. 请不要告诉我，这个人或那个人是逃跑的奴隶，或是强盗、小偷，或是充满无数过错，或是乞丐和卑贱，或是低贱、不值一提；但要想到，基督为他而死；这足以使你成为一切关怀的根基。要想想\Emph{他}是怎样的人，基督以如此高的代价估价他，甚至没有顾惜自己的血。因为，如果一位君王选择为任何人牺牲自己，我们不会再去寻找别的证明，证明此人是一位伟大而深得君王重视的人——我想不会——因为他的死亡足以显示死去的人对他的爱。然而，那不是人，不是天使，不是总领天使；而是诸天的主自己，天主的独生子自己，披上了血肉，为我们自由地交付了自己。我们难道不应做一切，尽一切努力，使那些被如此估价的人在我们手中享有完全的关怀吗？我们还将有什么辩护？什么宽容？\Person{保禄}在宣告这事时至少说了：\BibleQuote{不可因你的食物而毁坏基督为他而死的人。}\BibleRef{罗 14:15} 因为他愿意羞辱那些轻视弟兄、看不起他们软弱的人，并促使他们关怀邻人，说服他们照顾邻人，因此他放弃了其他一切，而提出了主的死亡。\par

那时，保禄坐在狱中，从那么遥远的地方写信给斐理伯人。因为按照天主的爱就是这样：任何属人的事物都不能打断它，因为它的根源来自天上，并有其赏报。他说什么呢？\Quote{弟兄们！我愿意你们知道。}\BibleRef{斐 1:12}你们看见他对学生的关怀吗？你们看见为人师者的细心吗？也请听听学生对老师的挚爱，好使你们知道，正是这种彼此相连的爱使他们坚强且不可战胜。因为如果\Quote{弟兄辅助弟兄，犹如坚固的城邑}，那么许多被爱的绳索捆绑在一起的人，更能完全击退恶魔的阴谋。的确，保禄与门徒相连，甚至不需要我们再加证明或辩论，因为事实上，他即使在捆绑中也殷切地关心他们，每天也为他们而死，心中燃烧着对他们的渴望。\par

6. 门徒们也是以完全的美德与保禄相连；不但是男人，连女人也是如此。关于福依贝，听他怎么说：\Quote{我把我们的姊妹福依贝推荐给你们，她是坚革哩教会的\OriginalTerm{女执事}{ministra}，你们要在主内以合乎圣徒的身份接待她，并在她需要你们的事上协助她，因为她帮助了许多人，也帮助了我自己。}\BibleRef{罗 16:1}但在这一事例中，他只就帮助方面为她作证；但普黎斯加和阿桂拉甚至为保禄舍命；关于他们，他这样写道：\Quote{阿桂拉和普黎斯加问候你们，他们为了我的性命，曾经冒过自己的颈项；}\BibleRef{罗 16:3}显然是至死。关于另一个人，在写给这些人的信中又说：\Quote{因为他为了基督的工作，曾冒死亡的危险，不顾自己的性命，为要补足你们在服事我上的缺欠。}\BibleRef{斐 2:30}你们看见他们如何爱自己的老师吗？他们如何把老师的安息看得比自己的生命还重要吗？因此，当时没有人能超越他们。我说这话，不只是为了让我们听见，更是为了效法；也不只是对属下说的，也是对治理者说的；好使学生们对老师表现出极大的关怀，老师们也像保禄一样，对属下怀有同样的挚爱；不只是对近处的，也对远处的。因为保禄以整个天下如一家，就这样不断地为所有人的得救操心；放下自己的事，捆绑、患难、鞭打和困境，每天查问门徒们的状况如何；并且常为此目的单独派人，时而派弟茂德，时而派提希苛；关于他，他说：\Quote{为叫你们知道我们的情形，并安慰你们的心；}\BibleRef{弗 6:22}关于弟茂德：\Quote{我打发了他，因为不能再忍，免得那诱惑者诱惑了你们。}\BibleRef{得前 3:5}又另处派了弟铎，又派了别人到别处。因为他自己因捆绑的束缚常被拘留在一处，不能与那些如同自己性命的人相会，便通过门徒们与他们相会。\par

7. 因此，那时他在捆绑中写信给斐理伯人，说：\Quote{弟兄们！我愿意你们知道，}\BibleRef{斐 1:12}称门徒为弟兄。因为爱就是这样：它消除一切不平等，不知有上下和尊卑；即使一个人高于众人，他也降到自己之下最卑微的位置；保禄正是这样做的。但让我们听听他愿意他们知道什么。\Quote{我所遭遇的事，更促进了福音的进展，}\BibleRef{斐 1:12}他说。告诉我，如何且以什么方式？难道你已脱了捆绑？解了锁链？能自由地在城里宣讲？难道你进入会堂，发表了长篇大论关于信仰的演讲，然后得了许多门徒离开？难道你复活了死人，成为众人惊叹的对象？洁净了癞病人，众人震惊？驱走了魔鬼，被高举？没有一件这样的事，他说。那么福音的进展如何发生？告诉我。\Quote{以致我的锁链在\OriginalTerm{整个宫廷}{toto praetorio}和其余众人中，已成了显明的事。}\BibleRef{斐 1:13}你说什么？这，这才是进展，这才是进步，这才是宣讲的增长——众人都知道你被捆绑了。是的，他说：至少听听下文，好叫你们知道，捆绑不但没有成为阻碍，反而成了更大的宣讲胆量的理由。\Quote{因此，许多在主内的弟兄，因我的锁链而信赖，更加大胆地敢宣讲天主的圣言。}\BibleRef{斐 1:14}保禄啊，你说什么？你的捆绑没有引起焦虑，反而引起信心？不是恐惧，而是热切的渴望？所说的事没有连贯性。我也知道。因为他说，这些事不是按照人间事理的连贯性发生的，而是所发生的事超越自然，成功来自天主的恩宠。因此，凡给别人带来焦虑的事，却给他带来了信心。因为如果有人擒获了军队的统帅并把他关起来，公之于众，他就使整个军营陷入溃逃；如果有人从羊群中劫走了牧人，他驱赶羊群的安全感就很大。但在保禄身上却不是这样，反而完全相反。因为军队的统帅被捆绑，士兵们却更加奋勇；他们向敌人冲击的信心更大：牧人被囚，羊群却没有被消灭，也没有被驱散。\par

8. 有谁曾见过，有谁曾听过，门徒们在师傅的危险中受到更大的鼓励？他们怎么不怕？他们怎么不惊恐？他们怎么没有对保禄说：\BibleQuote{医生，医治你自己吧！}\BibleRef{路 4:23} 拯救你自己脱离众多的危险，那时你才能为我们获得那无数好事？他们怎么没有说这些话？怎么！这是因为他们借圣神的恩宠受过训练，知道这些事发生并非由于软弱，而是出于基督的许可；为要使真理更广泛地照耀；借着锁链、监禁、磨难、困迫，不断增多升高，达到更大的规模。如此，基督的能力在软弱中得以成全。因为如果他的锁链使保禄残废，使他胆怯；无论对他自己还是对他的人；人不能不感到困难；但如果这些反而预备他获得更大的声望，人就必须惊愕和惊奇，如何借着涉及羞辱的事，为门徒赢得了光荣——借着引起胆怯的事，为他们众人带来了信心和鼓励。那时，谁看见他被锁链环绕，不感到惊愕？恶魔看见他在监狱里度日，更是纷纷逃遁。因为王冠使君王的头辉煌，不如锁链使他的双手辉煌；这不是由于它们自身的性质，而是由于那发出光明的恩宠。因此，门徒们得到了巨大的鼓励。因为他们看见他的身体虽然被捆绑，但他的舌头未被捆绑；他的手虽然被紧紧锁住，但他的声音未被束缚，并且比太阳光线更快地穿越整个世界。这成了对他们的鼓励，因为他们从事实中学到：没有一样现在的事是可畏惧的。因为当灵魂被神圣的渴求与爱真正浸透时，它不再关心任何现在的事；正如那些疯狂的人把自己投向火、剑、野兽、海洋和一切，这些人也一样，被一种最高贵、最属灵的狂热所疯狂，一种源于明智的狂热，嘲笑所有可见的事物。因此，看见他们的师傅被捆绑，他们更加欢欣，更加自豪；用事实向他们的对手证明他们无懈可击、不可征服。\par

9. 当下，当事情处于这种状态时，保禄的一些仇敌，想要把战争煽动得更加猛烈，并使暴君对他的仇恨加深，就假装他们自己也宣讲；（并且他们宣讲的是正确而纯正的信仰）为了教义更快地推进：他们这样做，并非出于传播信仰的渴望；而是为了让尼禄知道宣讲在增长、教义在推进，从而更快地把保禄处死。因此有两派：一派是保禄的门徒，一派是保禄的仇敌；一派出于真诚宣讲，另一派出于好斗和对保禄的仇恨。为表明这一点，他说：\BibleQuote{有的人固然出于嫉妒和纷争宣讲基督，但也有人出于善意；}（指出那些是他的仇敌）这是说他自己的门徒。\BibleRef{斐 1:15} 然后关于那些人：\BibleQuote{有的出于好争，}（他的仇敌）不纯洁，不纯正，而是\BibleQuote{以为这样可以给我的锁链带来压力；}\BibleRef{斐 5:17} \BibleQuote{但其他人出于爱；}（这又是关于他自己的弟兄）\BibleQuote{知道我是为辩护福音而设立的。}为何？然而无论如何，无论是假意还是真诚，基督总被传扬。\BibleRef{斐 1:16} 因此，这句引文若被理解为指异端，就是枉然无益。因为那时宣讲的人并不宣讲败坏的教义，而是宣讲正确纯正的信仰。如果他们宣讲败坏的教义，并教导与保禄相反的事，那么他们所渴望的必然不会成功。他们所渴望的是什么？那就是信仰增长，保禄的门徒增多，从而激起尼禄更大的敌意。如果他们宣讲不同的教义，就不会使保禄的门徒众多；如果门徒不众多，就不会激怒暴君。因此他并不是说他们引入败坏的教义——而是说他们宣讲的动机，\Emph{这一点}是败坏的。因为说明他们宣讲的借口本身不纯正，这是一回事。教义充满欺骗时，宣讲就不纯正；当宣讲确实纯正，但宣讲者不是为天主而宣讲，而是出于敌意或为了讨他人欢心时，借口就是不纯正的。\par

10. \Person{保禄}这样说，并不是因为他们引入了异端，而是因为他们传道并非出于正确的动机，也不是出于虔诚。因为他们这样做不是为了扩展福音，而是为了与他争战，将他置于更大的危险中——正是为此他指责他们。请看何等精确地陈述。他说：\Quote{以为，}他说，\Quote{他们是给捆锁我的人加压力。}\BibleRef{斐 1:17}他没有说\Quote{施加}，而是说\Quote{以为他们是施加}，即假设，以此指明即使他们这样假设，他自己仍不处于那种境地；反而因福音的进展而喜乐。因此他补充说：\Quote{但在这事上，我既已喜乐，也还要喜乐。}\BibleRef{斐 1:18}倘若他认为他们的教义是欺骗，且他们引入了异端，那么\Person{保禄}不可能喜乐。但由于教义是纯正且出自真正源头，因此他说：\Quote{我喜乐，也还要喜乐。}因为他们若因好争而毁灭自己，那又如何呢？即便如此，他们不情愿地也增强了我事业的力量。你们看到\Person{保禄}的能力何等伟大吗？他如何不被魔鬼的任何诡计所捕获？他非但未被捕获，反而藉这些自身制服了魔鬼。因为魔鬼的狡猾和那些服事它之人的邪恶确实极大；他们假借同心合意的外表，想要消灭这宣讲。但\Quote{那在狡猾者的诡计中抓住他们的}并不允许那时发生这事。至少\Person{保禄}在宣告此事时说：\Quote{但我在肉身继续存留，为你们更是必要的；我确切知道这一点，即我将继续并留在你们众人中间。}\BibleRef{斐 1:24}那些人确实一心要将我从今生逐出，并为此预备忍受一切；但天主为了你们的缘故不允许这事。\par

11. 因此，你们要精确地记住这一切，以便能运用全备的智慧纠正那些随意使用圣经而无视上下文、且为毁灭邻人的人。我们要能记住所说的，并纠正他人，就必须常以祈祷为避难所，恳求那赐予智慧之言的天主，赐给我们聆听的悟性，并对交在我们手中的这属灵宝库予以谨慎而不可征服的守护。因为许多事我们靠自己努力无力成功完成，却可以通过坚持不懈的祈祷轻易达成。因为无论处于患难、危险或顺境中——处于安逸和极大昌盛中的人，要祈祷这些保持不变、没有变迁、永不改变；处于患难和诸多危险中的人，要祈祷看到某种有利的变化临到他，并被带入慰藉的平静中——时常不间断地祈祷是当行之责。你处于平静中吗？那么恳求天主使这平静继续稳固地在你身上。你看见风暴兴起反对你吗？热切恳求天主使波涛过去，并从中得着平静。你已被垂听了吗？要为此衷心感恩，因你已被垂听。你未被垂听吗？坚持，以便你被垂听。因为即使天主有时迟延给予，也不是出于仇恨和厌恶；而是想通过迟延给予，使你永远留在祂身边，正如慈爱的父亲们所做的那样：他们通过迟延给予，巧妙地管理那些较为懒惰的孩子们的持久和殷勤的陪伴。你在与天主会面时，无需中介，无需太多拉拢，无需谄媚他人；即使你孤苦伶仃，即使没有代言人，独自一人，呼求天主帮助，你定能成功。祂不轻易应允别人为我们代求，却应允我们这些有需要者自己祈祷，即使我们背负着万千恶行。因为与人相处时，即使我们与他们发生无数冲突，当我们在清晨、中午和晚上不断出现在那些对我们怀怨的人面前时，通过不间断的坚持和持续的会面交谈，我们便轻易消除了他们的敌意——何况在天主面前，更是如此。\par

12. 但你不配。凭你的\OriginalTerm{恒心}{assiduity}成为配得吧。因为以下三点：不配的人凭其恒心成为配得是可能的；我们亲自呼求天主比藉他人呼求更能得到主的同意；主常常延迟施予，并非希望我们完全困惑，也不是要我们空手而去，而是为使我们成为更大恩惠的领受者——我将努力借今天读给你们的那段比喻来阐明这三点。一位客纳罕妇人来到基督面前，为一个被魔鬼附身的女儿祈求，极其恳切地呼喊着（经文说：\BibleQuote{主啊，可怜我吧，我的女儿被魔鬼纠缠得苦。}）看，这妇人来自异族，是外邦人，在犹太社团之外。她实在与一条狗无异，不配获得所求。因为他说：\BibleQuote{拿儿女的饼扔给小狗是不对的。}然而，正是因她的恒心，她成了配得的。因为他不仅使她这曾为狗的，得以跻身儿女的尊贵之列，还以那崇高的赞辞送她离去，说：\BibleQuote{妇人，你的信德真大！照你所愿，给你成就吧。}当基督说\BibleQuote{你的信德真大}时，你无需再寻求其他证据来证明这妇人心灵的伟大。你看见了吗？这妇人如何因她的恒心，从不配成为配得？你岂不愿意也学习：我们亲自呼求他比藉他人呼求更能成就我们的愿望？她大声呼喊，门徒们来到他跟前说：\BibleQuote{请打发她走吧，因为她在我们后面喊叫。}\BibleRef{玛 5:23} 他却对他们说：\BibleQuote{我被派遣，只是为了以色列家迷失的羊。}\BibleRef{玛 5:24}但是当她亲自来到他面前，继续呼喊说：\BibleQuote{主啊，是的，可是小狗也吃主人桌子上掉下来的碎屑。}那时他应允了她的请求，说：\BibleQuote{照你所愿，给你成就吧。}你看见了吗？当他们为他求情时，他拒绝了；而当那需要恩赐的妇人亲自呼喊时，他同意了？因为对他们，他说：\BibleQuote{我被派遣，只是为了以色列家迷失的羊；}但对她说：\BibleQuote{你的信德真大！照你所愿，给你成就吧。}再者，在她祈求的开始和序奏中，他什么也没有回答；但当她一次、两次、三次来到他面前时，他才施予恩赐；藉此结局使我们相信，他延迟施予，并非要拒绝她，而是要向我们彰显这妇人全部的忍耐。因为如果他延迟是为了拒绝她，他最终也不会应允；但他等待是为了向所有人显明她的神性智慧，因此他保持沉默。因为如果他立刻并在开始时应允，我们就不会认识这妇人的美德。经上说：\BibleQuote{请打发她走吧，因为她在我们后面叫喊。}但基督说什么？你们听见声音，我却看见心思：我知道她将说什么。我选择不让隐藏在她心里的宝藏被忽视；我等待并保持沉默，为的是在发现它之后，公之于众，并向所有人显明。\par

13. 因此，既然学习了这一切，即使我们处于罪中，不配领受，也不要绝望；知道凭着心灵的恒心，我们将能成为配得所求的人。即使没有代祷者帮助，处于贫乏之中，也不要气馁；知道亲自满怀热忱来到天主面前，本身就是强有力的代祷。即使他延迟施予，也不要灰心；已经明白拖延和延迟正是他眷顾和慈爱的确实证据。如果我们这样说服自己；带着深痛而热切的灵魂、全然激发的决心，如同那位客纳罕妇人前来时那样，我们也来到他面前，即使我们是狗；即使我们做过任何可怕的事；我们既能驳斥自己的罪行，也能获得如此大的放胆直言，以致还能为他人代祷；正如这位客纳罕妇人，她不仅自己享有放胆直言和无数赞美，还有能力将她亲爱的女儿从难忍的痛苦中救出。因为没有——没有比热诚而真诚的祈祷更有力的了。这既驱散眼前的危险，也从那时刻的刑罚中拯救出来。因此，为使我们能安然度过今生，并满怀信心离世，让我们以极大的热忱和渴望恒常实行祈祷。因为这样我们就能获得那为我们存留的美善，并享有那些卓越的盼望；愿天主恩赐我们众人获得；藉着我们的主耶稣基督的恩宠、慈爱和怜悯——愿光荣、尊威、权柄归于他，及圣父与圣神，至于无穷世之世。阿们。\par

\SourceRecord{Translated by R. Blackburn.；Nicene and Post-Nicene Fathers, First Series；Vol. 9.；Edited by Philip Schaff.；Buffalo, NY: Christian Literature Publishing Co.,；1889.；<http://www.newadvent.org/fathers/1907.htm>.}
